\subsection{HOMOGENEOUS SPACES}

Let G be a topological group and H a subgroup of G. H operates continuously on the right on G according to the external law \((t, x) \to xt\) , and the orbit of a point \(x \in G\) is the left coset xH. The set of orbits is therefore what we have called in algebra the homogeneous space G/H. Whenever we speak of G/H as a topological space, we shall always mean the orbit space of G (with respect to H) unless the contrary is expressly stated; i.e.~the quotient space of G by the equivalence relation \(x^{-1}y \in H\) . In conformity with the general definitions, we say that the topology of this space is the quotient by H of the topology of G.

PROPOSITION 12. The group G operates continuously on every homogeneous space G/H.

Since the equivalence relation \(x^{-1}y \in H\) is open (no. 4, Lemma 2) this is a particular case of Proposition 11 of no. 4.

PROPOSITION 13. Let G be a topological group and let H be a subgroup of G. Then the homogeneous space G/H is Hausdorff if and only if H is closed in G.

H is an equivalence class for the relation \(x^{-1}y \in H\) and therefore, if G/H is Hausdorff, H is closed in G. Conversely, if H is closed, then the graph of this relation is closed in \(G \times G\) , since it is the inverse image of H under the continuous mapping \((x, y) \to x^{-1}y\) . Since the relation \(x^{-1}y \in H\) is open, it follows from Chapter I, § 8, no. 3, Proposition 8 that G/H is Hausdorff.

PROPOSITION 14. Let G be a topological group and let H be a subgroup of G. Then the homogeneous space G/H is discrete if and only if H is open in G.

For the inverse images in G of the points of G/H under the canonical mapping are the cosets xH (x ∈ G); and these sets are open in G if and only if H is open in G.

Let X be a topological space on which a topological group G operates continuously and transitively; X is then (in the algebraic sense) a homogeneous space of G. Let \(x\) be a point of \(X\), \(H_x\) its stabilizer. The continuous surjection \(s \to s.x\) of \(G\) onto \(X\) factorizes canonically as

\[
\mathbf {G} \xrightarrow {f _ {x}} \mathbf {G} / \mathbf {H} _ {x} \xrightarrow {g _ {x}} \mathbf {X}
\]

\pandocbounded{\includegraphics[keepaspectratio]{images/part002_c9cfa224.jpg}}

where \(f_x\) is the canonical mapping of \(G\) onto the homogeneous space \(G/H_x\) , and \(g_x\) is the bijection \(s.H_x \to s.x\) of \(G/H_x\) onto \(X\) ; moreover (Chapter I, § 3, no. 4, Proposition 6) \(g_x\) is a continuous mapping. But \(g_x\) is not necessarily a homeomorphism of \(G/H_x\) onto \(X\) (Exercise 29). If \(g_x\) is a homeomorphism for each \(x \in X\) , then we say that \(X\) is a topological homogeneous space (of the topological group \(G\) ); for this to be so it is necessary and sufficient that, for each \(x \in E\) , the mapping \(s \to s.x\) of \(G\) into \(X\) should be open.

PROPOSITION 15. Let X be a topological space on which a topological group G operates continuously and transitively. For X to be a topological homogeneous space (relative to G) it is sufficient that for some point \(x_{0} \in X\) the mapping \(s \to s \cdot x_{0}\) transforms each neighbourhood of e in G into a neighbourhood of \(x_{0}\) in X.

Every \(x \in X\) can be written as \(x = t.x_0\) for some \(t \in G\) . If \(V\) is a neighbourhood of \(e\) , then \(V.x = (Vt).x_0\) is a neighbourhood of \(x\) , for we can write \((Vt).x_0 = t((t^{-1}Vt).x_0)\) , and the assertion follows from the facts that \(t^{-1}Vt\) is a neighbourhood of \(e\) in \(G\) and that \(y \to t.y\) is a homeomorphism of \(X\) onto itself (no. 4, Lemma 1). It follows that if \(U\) is any open subset of \(G\) and if \(x\) is any point of \(X\) , then \(U.x\) is open in \(X\) ; for if \(t \in U\) , then \(t^{-1}U\) is a neighbourhood of \(e\) , hence \((t^{-1}U).x\) is a neighbourhood of \(x\) , and \(t.((t^{-1}U).x) = U.x\) is a neighbourhood of \(t.x\) . Hence \(U.x\) is open in \(X\) , so that the mapping \(s \to s.x\) of \(G\) into \(X\) is open. This completes the proof.

